\scanpage{30}
\mspage{6 - 41 -}
n\textsuperscript{o} 392

\S~5.

1) Soient $A$ un anneau noethérien intégralement clos, $K$ son corps des fractions, $L$ une extension de degré $n$ de $K$, $B$ la fermeture intégrale de $A$ dans $L$ ; on suppose que $B$ est un $A$-module de type fini. Soient $E$, $F$ deux anneaux tels que $A\subset E\subset F\subset B$.

a) Montrer que pour tout idéal premier $\underline p$ de hauteur $1$ dans $E$, il n'y a qu'un nombre fini d'idéaux premiers $\underline P$ de $F$ au-dessus de $\underline p$ et qu'ils sont de hauteur $1$ ; réciproquement, pour tout idéal premier $\underline P$ de hauteur $1$ dans $F$, $\underline P\cap E$ est de hauteur $1$ dans $E$.

b) Pour tout idéal premier $\underline P$ de hauteur $1$ dans $F$, on appelle \uline{degré} de $\underline P$ par rapport à $E$, et on note $\deg_E(\underline P)$ le degré du corps des fractions de $F/\underline P$ sur le corps des fractions de $E/(\underline P\cap E)$ ; le degré par rapport à $E$ d'un idéal premier de $F$ qui n'est pas de hauteur $1$ est pris égal à\unsure{.|,} $0$ par définition. Si $\underline Q$ est un idéal primaire de $F$, $\underline P$ l'idéal premier correspondant, et si $m$ est le plus petit entier tel que la puissance symbolique $\underline P^{(m)}$ de $P$ soit contenue dans $\underline Q$, on pose $\deg_E(\underline Q)=m.\deg_E(\underline P)$, (de sorte que $\deg_F(\underline Q)=m$ si $\underline P$ est de hauteur $1$, $\deg_F(\underline Q)=0$ dans le cas contraire). Finalement, si $\underline a$ est un idéal quelconque $\ne 0$ de $F$, $\underline Q_i$ les idéaux primaires de la décomposition primaire réduite de $\underline a$, correspondant aux idéaux premiers non immergés, on pose $\deg_E(\underline a)=\sum_i\deg_E(\underline Q_i)$.

Avec ces notations, soit $C$ un anneau tel que $A\subset C\subset B$ ; on dit qu'un idéal $\underline a$ de $C$ est \uline{pseudo-inversible} si, pour tout idéal premier $\underline p$ de hauteur $1$ dans $A$, $\underline a_{\underline p}$ est inversible dans l'anneau semi-local $C_{\underline p}$. Montrer que, pour que $\underline a\ne 0$ soit pseudo-inversible, il faut et il suffit que l'on ait $\deg_C(\underline a)=\sum_{\underline P}n(\underline P)\deg_C(\underline P)$, o\rlap{\handwritten{\`{\phantom{u}}}}u $\underline P$ parcourt l'ensemble des idéaux premiers de hauteur $1$ de $B$ et les $n(\underline P)$ sont définis par $\operatorname{div}(\underline a B)=\sum_{\underline P}n(\underline P)\underline P$. (Se ramener au cas o\rlap{\handwritten{\`{\phantom{u}}}}u $A$ est un anneau de valuation discrète, donc $C$ semi-local. Observer d'une part que dans $C$, il y a alors identité entre les idéaux inversibles et les idéaux principaux, et que si $\underline a=xC$, $\deg_C(\underline a)$ est égal à la longueur du $C$-module $C/xC$, et aussi à celle du $C$-module $B/xB$. D'autre part, si $\underline m_i$, $(1\leq i\leq r)$, sont les idéaux maximaux distincts de $C$, $\underline a\ne 0$ un idéal de $C$, il existe $y\in\underline a$ tel que $y\notin C\cap(\underline a\underline m_i B)$ pour tout $i$ ; en déduire que $yB=\underline a B$, en utilisant le fait que $B$ est alors un anneau principal). Si $C=B$, tout idéal $\ne 0$ de $B$ est pseudo-inversible.
% archive sheet number: 29

